% Imporditud: tools/import_eksperimendid.py
% Allikas: Eesti füüsikaolümpiaadi eksperimendi ülesannete
%   kogu (Rael Kalda, 2023), ülesanne Ü167, lahendus L167.
\setAuthor{EFO žürii}
\setRound{lõppvoor}
\setYear{2008}
\setNumber{G E2}
\setDifficulty{5}
\setTopic{dünaamika}
\setVahendid{Kolm lehte paberit, käärid, dünamomeeter (mõõtepiirkond 5 N), koormis (massiga 65 $\pm$ 0,5 g)}
\setPunktid{14}
\setAllikas{Eksperimendi kogu Ü167, lahendus lk 125}
\setLahendusseis{toores}

\prob{Paber}
Määrake hõõrdetegur paberi ja paberi vahel parima võimaliku täpsusega. Hinnake mõõteviga.

\hint

\solu
Asetame kaks paberilehte teineteise peale, kinnitame ühe külge dünamomeetri, koormame ühenduskoha koormisega ning mõõdame lehtede vahel tekkinud hõõrdejõu. Saame tulemuseks

Fh = 0,2 $\pm$ 0,1 N,

(kus mõõteriista viga koosneb nii nullinihkest kui ka süstemaatilisest veast). Mõõtmistulemus on ilmselgelt mõõtevea piiril ja ei ole ligilähedaseltki sama suhtelise veaga läbi viidud mõõtmine kui massi mõõtmise suhteline viga. Mõõtmistäpsuse suurendamiseks on vaja suurendada hõõrdejõudu. Rõhumisjõudu suurendada pole võimalik. Küll on aga võimalik suurendada samaaegselt kokku puutuvate ühe ja sama rõhumisjõu all olevate hõõrduvate pindade arvu. Summaarne saadav hõõrdejõud oleks siis

FSh = Fh $\cdot$ N,

kus N on samaaegselt kokku puutuvate hõõrduvate pindade arv. Lõikame paberist ribad ja asetame neid suurema hulga vaheldumisi osaliselt teineteisega üle kattuma nii, et üle ühe oleks võimalik lehti kinni hoida ja vahepealsed lehed oleksid üheagselt kinnitatud dünamomeetri külge. Asetame vaheldumisi 11 riba saades kokkupuutuvate pindade arvuks 10 ning koormame meie koormisega ning mõõdame summaarse hõõrdejõu. Tulemus

FSh = 2,3 $\pm$ 0,1 N on 10 korda täpsem kui esialgselt mõõdetud. Lõppavaldis on

$\mu$ = FSh , Nmg

arvuliselt $\mu$ = 0,35 $\pm$ 0,02.
\probend
